\chapter{Pediatric Imaging}

\label{ch:ch56_pediatric}

\section{Why is pediatric imaging different?}
Children are not small adults. Their anatomy is developing, their disease spectrum is different, and they are more sensitive to ionising radiation. Pediatric imaging demands a thoughtful approach that balances diagnostic information against long-term radiation risk, addresses the challenges of imaging uncooperative patients, and recognises age-specific normal findings that could be misinterpreted as pathology.

\section{ALARA: radiation dose optimisation}
The ALARA principle (As Low As Reasonably Achievable) is foundational in pediatric imaging:
\begin{itemize}
    \item \textbf{Use ultrasound or MRI when possible} to avoid radiation entirely.
    \item \textbf{When CT is necessary, use pediatric protocols} with reduced tube current and voltage.
    \item \textbf{Limit the scan range} to the area of clinical interest only.
    \item \textbf{Use iterative reconstruction} to reduce noise at lower doses.
    \item \textbf{Single-phase imaging} where multiphasic studies are not required.
    \item \textbf{Document dose length product (DLP)} and compare with local diagnostic reference levels.
\end{itemize}

The Image Gently campaign promotes radiation safety in children globally. Its core message: ``Go low, go slow, and say no'' to unnecessary CT.

\begin{pearlbox}[CT dose in children]
A single abdominal CT in a child delivers an effective dose equivalent to 1--5 years of background radiation. The cancer risk from pediatric CT is estimated at 1 excess case per 1000--10,000 scans, depending on age and organ dose. This small but real risk must be weighed against the diagnostic benefit.
\end{pearlbox}

\section{Imaging in neonates}

\subsection{Neonatal head ultrasound}
Head ultrasound is the first-line imaging investigation in neonates (particularly preterm infants):
\begin{itemize}
    \item \textbf{Indications:} Screening for intraventricular hemorrhage (IVH), periventricular leukomalacia (PVL), hydrocephalus, and congenital anomalies.
    \item \textbf{Technique:} Performed through the anterior fontanelle as an acoustic window. No sedation or radiation required.
    \item \textbf{Normal findings:} The germinal matrix is visible as an echogenic area in the caudothalamic groove and involutes over the first weeks. The cavum septum pellucidum and cavum Vergae are normal findings.
    \item \textbf{IVH grading:}
\end{itemize}

\begin{table}[h]
\centering
\caption{Grading of intraventricular hemorrhage}
\label{tab:ivh}
\begin{tabular}{cll}
\hline
\textbf{Grade} & \textbf{Description} & \textbf{Prognosis} \\
\hline
I & Germinal matrix hemorrhage only & Excellent \\
II & IVH without ventricular dilatation & Good \\
III & IVH with ventricular dilatation & Guarded (risk of PVL) \\
IV & IVH with parenchymal involvement & Poor (cerebral palsy risk) \\
\hline
\end{tabular}
\end{table}

\subsection{Hip ultrasound in infants}
Developmental dysplasia of the hip (DDH) is screened with ultrasound in the first 6 weeks of life:
\begin{itemize}
    \item \textbf{Graf method:} Classifies the hip by alpha angle (bony roof angle) and beta angle (cartilaginous roof angle).
    \item \textbf{Normal (Type I):} Alpha angle $>$60$^\circ$, femoral head well-covered.
    \item \textbf{Abnormal (Type II--V):} Decreased alpha angle, uncovering of femoral head. Types III--V require treatment (Pavlik harness).
    \item \textbf{Indications:} Clinical examination abnormality, breech presentation, family history, oligohydramnios.
    \item \textbf{After 4--6 months,} ossification of the femoral head limits ultrasound utility; X-ray (frog-leg lateral) is used instead.
\end{itemize}

\section{Pediatric fractures}

\subsection{Growth plate injuries}
Growth plate (physeal) injuries are unique to children and classified by the Salter-Harris system (see Chapter~\ref{ch:ch54_musculoskeletal}). Key points:
\begin{itemize}
    \item The physis is the weakest link in the paedicular bone-ligament unit---children get physeal fractures where adults get ligament tears.
    \item Types I and II have excellent prognosis. Types III, IV, and V carry risk of growth arrest.
    \item A suspected Type I fracture with normal X-ray should be treated clinically and splinted; repeat X-ray in 10--14 days or MRI if concern persists.
\end{itemize}

\subsection{Toddler's fracture}
A toddler's fracture is a spiral or oblique fracture of the distal tibial shaft in children aged 9 months to 3 years, typically caused by a twisting injury.
\begin{itemize}
    \item Often clinically subtle---the child refuses to bear weight.
    \item The fracture may be invisible on initial X-ray.
    \item If clinical suspicion remains high despite a normal X-ray, treat with a backslab and repeat X-ray in 10--14 days. Healing callus will then confirm the diagnosis.
\end{itemize}

\begin{pearlbox}[Occult fractures in children]
If a child has focal tenderness and refusal to use a limb, but X-rays are normal, do not dismiss the presentation. Treat as a fracture, arrange splinting, and repeat imaging in 10--14 days. MRI is the most sensitive modality for occult fractures and is preferable to repeat CT.
\end{pearlbox}

\section{Pediatric chest imaging}

\subsection{Croup (laryngotracheobronchitis)}
Croup is a common viral illness causing upper airway obstruction in children aged 6 months to 3 years.
\begin{itemize}
    \item \textbf{Clinical:} Barking cough, stridor, hoarseness.
    \item \textbf{Lateral neck X-ray:} ``Steeple sign''---symmetric narrowing of the subglottic trachea. This finding is present in approximately 50\% of cases.
    \item \textbf{Usually clinical, not radiological:} Imaging is only required when the diagnosis is uncertain or to exclude other causes of stridor.
\end{itemize}

\subsection{Epiglottitis}
Epiglottitis is an emergency caused by infection of the epiglottis, potentially leading to complete airway obstruction.
\begin{itemize}
    \item \textbf{Clinical:} Sudden onset fever, drooling, dysphagia, toxic appearance, tripod positioning.
    \item \textbf{Lateral neck X-ray:} ``Thumb sign''---swollen, rounded epiglottis. Thickened aryepiglottic folds.
    \item \textbf{Important:} Do \emph{not} attempt to examine the throat with a tongue depressor---this can precipitate complete airway obstruction. Secure the airway first.
\end{itemize}

\section{Appendicitis in children}
Appendicitis is the most common surgical emergency in children.
\begin{itemize}
    \item \textbf{Ultrasound (first-line):} Non-compressible, aperistaltic, blind-ending tubular structure $>$6~mm in diameter, with appendicolith in 30\% of cases. Periappendiceal fluid and increased vascularity on Doppler support the diagnosis.
    \item \textbf{CT (second-line):} Used when ultrasound is equivocal. Findings include a dilated appendix ($>$6~mm), wall enhancement, periappendiceal fat stranding, appendicolith, and abscess.
    \item \textbf{MRI (alternative):} Increasingly used as a radiation-free alternative in children, particularly in centers with rapid MRI protocols.
    \item \textbf{Dose consideration:} CT for appendicitis in children delivers a significant abdominal dose. Use MRI or ultrasound as the primary modality wherever possible.
\end{itemize}

\section{Imaging in child abuse (non-accidental injury)}
Imaging plays a critical role in the evaluation of suspected child abuse:
\begin{itemize}
    \item \textbf{Skeletal survey:} A complete radiographic skeletal survey (minimum 21 views) is indicated in all children $<$2 years with suspected non-accidental injury (NAI). Look for:
    \begin{itemize}
        \item Metaphyseal corner fractures (``bucket handle'' or ``corner'' fractures)---highly specific for NAI
        \item Rib fractures (especially posterior---suggestive of squeezing injury)
        \item Fractures in different stages of healing
        \item Spiral fractures of long bones in non-ambulant children
    \end{itemize}
    \item \textbf{CT head:} Indicated if intracranial injury is suspected (altered consciousness, seizures, vomiting). Look for subdural hemorrhage (bridging vein rupture), retinal hemorrhage (ophthalmology), and cerebral edema.
    \item \textbf{Abdominal imaging:} CT abdomen/pelvis if there is concern for solid organ injury (abdominal tenderness, bruising, haematuria).
\end{itemize}

\begin{pearlbox}[The radiographer's role]
Radiographers are often the first professionals to recognize patterns of injury inconsistent with the given history. A skeletal survey showing fractures of different ages in a child with a history of ``falling off the bed'' should raise immediate concern. Document findings carefully and follow local safeguarding protocols.
\end{pearlbox}

\section{Age-specific normal findings}
Knowledge of normal developmental anatomy is essential to avoid misdiagnosis:

\begin{table}[h]
\centering
\caption{Normal findings that may mimic pathology}
\label{tab:paediatric_normals}
\begin{tabular}{@{}p{0.33\textwidth}p{0.15\textwidth}p{0.39\textwidth}@{}}
\hline
\textbf{Finding} & \textbf{Age} & \textbf{Key point} \\
\hline
Nasopharyngeal lymphoid tissue & Toddlers & Can simulate nasopharyngeal mass \\
Ossification centers (wrist) & Variable & Follow age-appropriate charts \\
Physiological subperiosteal new bone & 1--4 months & Normal; resolves; do not mistake for abuse \\
Thymus (chest X-ray) & $<$2 years & Large thymus is normal; do not mistake for mass \\
Tibial tubercle prominence & 10--14 years & Normal variant; not a fracture \\
Tracheal indentation by aortic arch & Children & May appear as ``narrowing'' on AP view \\
\hline
\end{tabular}
\end{table}

\section{Pediatric sedation for imaging}
Some pediatric imaging requires sedation to prevent motion artefact:
\begin{itemize}
    \item \textbf{Indications:} MRI in children $<$6 years (sometimes up to 8 years), complex CT requiring motion-free images.
    \item \textbf{Options:} Chloral hydrate (oral), midazolam (intranasal/buccal), dexmedetomidine (intranasal), general anaesthesia (for lengthy procedures).
    \item \textbf{Alternatives to sedation:} Feed-and-sleep technique for MRI in infants, ``swaddling'' and immobilisation devices, child-life specialist distraction, rapid MRI protocols.
    \item \textbf{Safety:} All sedated children require monitoring (pulse oximetry, heart rate, respiratory rate) and must be recovered under observation. Follow local sedation policy and training requirements.
\end{itemize}

\begin{pearlbox}[Feed and sleep MRI]
For infants $<$6 months, MRI can often be performed without sedation by feeding and swaddling the infant immediately before the scan. This technique avoids the risks of sedation and is successful in 80--90\% of cases with appropriate preparation.
\end{pearlbox}

\section{When to escalate imaging}
Escalate from ultrasound or X-ray to CT/MRI when:
\begin{itemize}
    \item A suspicious bone lesion requires characterisation
    \item Occult fracture is suspected but X-ray is negative
    \item Suspected intracranial injury (abnormal neurology, altered consciousness)
    \item Complex appendicitis or intra-abdominal pathology
    \item Staging of a known malignancy
    \item Evaluation of a suspected congenital anomaly (MRI for CNS, CT for complex skeletal anomalies)
\end{itemize}

\section{Reporting checklist}
A structured pediatric report should include:
\begin{enumerate}
    \item Age-appropriate technique and dose optimisation
    \item Age-appropriate normal variants vs true pathology
    \item Growth plate assessment (if applicable)
    \item Fracture description (Salter-Harris type, displacement, angulation)
    \item Organ-specific findings (age-appropriate criteria)
    \item Comparison with prior imaging (growth and development)
    \item Clear recommendation for follow-up or escalation
    \item Safeguarding concerns (if applicable): follow local policy
\end{enumerate}

\clearpage
